% ============================================================================
\clearpage
\section*{v4.10.0-r08 Source-Only: The Assumptions of the Companion and the Dependence of the Open Predictions}
\addcontentsline{toc}{section}{v4.10.0-r08 Source-Only: The Assumptions of the Companion and the Dependence of the Open Predictions}
\label{sec:v4100-r08-assumptions}
% ============================================================================

\paragraph{Purpose.}
Dorau and Much (\S\ref{sec:v4100-r06-relative-entropy}) start from few assumptions, set openly and motivated independently, and derive one
strict consequence.  This layer applies that way of working to the companion.  First it settles a suspected conflict between the second reading
and the photon weights of gravity.  Then it lists the assumptions of the companion with their origin, and it traces in QAM bookkeeping what each of
the four open predictions rests on.  The layer is an inventory; it contains no test and no new prediction.  The notes are
\texttt{notes/\allowbreak research/\allowbreak v4\_10\_0\_r08/}.  Earlier layers are not rewritten.

\noindent\textbf{Source-only status.}
This v4.10.0-r08 layer is source-only.  It does not build a PDF, does not assign a new DOI,
and does not authorize a public release.  The v4.10.0 Version DOI
\href{https://doi.org/10.5281/zenodo.23991357}{10.5281/zenodo.23991357} and the Concept DOI
\href{https://doi.org/10.5281/zenodo.19210728}{10.5281/zenodo.19210728} are unchanged by it.

\subsection*{v4.10.0-r08 results}
\label{subsec:v4100-r08-results}

\paragraph{Theoremlet v4.10.0-r08.1 (the source of gravity under the weights).}
Exploration, analytic.  The suspicion was: under the second reading the action is \((a-b)S_M\) (H79), and with photon weights \(a=b\) for gravity
(v4.10.0-r05.1(b)) the gravitating source of matter would vanish.
\begin{enumerate}[leftmargin=2em]
\item[(a)] The sum \((a-b)\) arises only where both readings live at the same place (domain G1, explored in v4.10.0-r04, not adopted).  Under the
adopted domain G2 the matter action is \(a\,S_M\) after the bang and \(-b\,S_M\) on its image before; no sum occurs.
\item[(b)] The photon weights concern the gravitational action: \(-b_g S_{\rm grav}\) before, \(a_g S_{\rm grav}\) after the bang, \(a_g=b_g\).  After the bang
\(a_g\tilde G=a\,\kappa\,T_M\), so \(G_{\rm eff}=(a/a_g)\,G\), one common factor absorbed into the measured \(G\).  There is no conflict.
\item[(c)] \(G_{\rm eff}\) is universal only if all non-gravitational fields carry the same weight \(a\) after the bang.  Different overall factors for two
species would give them different active gravitational masses at equal inertial mass.  This assumption (U) is used tacitly in v4.10.0-r03 and is
stated nowhere; it is recorded here.  The factors of H82 act within one action and change parameters, not universality.
\item[(d)] The limit of r07 for massive fields concerns only the domain before the bang: under G2 only the matter reading lives after it, so the chain
of r06.1 applies today in frame M to massive and massless fields alike.
\end{enumerate}

\paragraph{Remark v4.10.0-r08.2 (the list of assumptions).}
The assumptions are listed in \texttt{Annahmenliste\_\allowbreak Companion.md} with origin and one of six categories: \emph{V} set before the consequences
were known, with a reason of its own; \emph{R} set afterwards, a selection or repair after a consequence was known or did not fit; \emph{K} correction of an
error; \emph{A} working assumption; \emph{L} input from the literature; \emph{S} tacit.  An R assumption is not a flaw, but a prediction resting on it is less
blind by that selection.  Of the action programme W1/W2 (frozen since v4.10.0-r05), seven entries are R or A within two days; this is the repair series
recorded in r05.2.  None of them enters the open predictions.

\paragraph{Theoremlet v4.10.0-r08.3 (what the open predictions rest on).}
Bookkeeping in QAM (\texttt{qam\_\allowbreak annahmen\_\allowbreak vorhersagen.tex}; \texttt{qam.sty} raises no error).  All four open predictions have status
\emph{reading}: their weakest premises are readings of the core, not derivations.  Their inputs according to the exported record:
\begin{center}
\small
\begin{tabular}{@{}>{\raggedright\arraybackslash}p{0.28\textwidth}>{\raggedright\arraybackslash}p{0.64\textwidth}@{}}
\toprule
prediction & inputs\\
\midrule
H96, \(\Omega_M=\beta_{\rm sym}^2\) & \(\gamma_{\rm BI}\), \(\gamma_{\rm string}\), counting rule, epoch ``today''\\
H15, \(m_W\) & \(\gamma_{\rm BI}\), \(\gamma_{\rm string}\), counting rule, \(m_Z\), reading \(\cot\theta_W=\gamma_L\)\\
H46, \(\sin^2\theta_{23}\in[0.49,0.50]\) & neutrino reading, mask, neutrino data, \(\varepsilon=\eta\), \(\gamma_{\rm BI}\), \(\gamma_{\rm string}\), counting rule\\
H42, \(\sin^2\theta_{23}\le0.54\) & neutrino reading, mask, neutrino data\\
\bottomrule
\end{tabular}
\end{center}
\begin{enumerate}[leftmargin=2em]
\item[(a)] \emph{H15, H46 and H96 rest on the same core}: the two literature values, the assignment of counting rules to the frames (M all spins, T
\(j=\tfrac12\), U(1) rather than SU(2); v4.9.0-r63, r68) and, through the frame map, the reading of \(\gamma_L\) as Lorentz factor.  If the counting rule
fails, all three fail together.  With the SU(2) counting preferred in the literature (\(\gamma_L\approx2.15\)) one would get \(m_W=82.68\) GeV instead of
\(80.361\) GeV and \(\Omega_M=0.365\) instead of \(0.302\).  The U(1) rule was selected by the lepton window (r63) and the \(m_W\) relation was noticed
afterwards, so these two agreements are not independent of the selection.
\item[(b)] H42 and H46 share the neutrino reading (r55, r57); H46 adds \(\varepsilon=\eta\) and with it the \(\gamma_L\) core.
\item[(c)] Only H42 is independent of the \(\gamma_L\) core.  The four open predictions form two independent groups, not four.
\item[(d)] Forming \(\gamma_L\) without the declared map from T to M is reported by QAM as frame mixing.  \(\gamma_L\) is defined only through the reading it
is used to support.
\end{enumerate}

\paragraph{Proof and measurement sources.}
\texttt{Erkundung\_\allowbreak Quelle\_\allowbreak Gravitation\_\allowbreak Gewichte.md} (r08.1); \texttt{Annahmenliste\_\allowbreak Companion.md} (r08.2, r08.3);
\texttt{qam\_\allowbreak annahmen\_\allowbreak H96.tex} and \texttt{qam\_\allowbreak annahmen\_\allowbreak vorhersagen.tex} with exported records.  The SU(2) values in r08.3(a)
follow from \(m_W=m_Z\gamma_L/\sqrt{1+\gamma_L^2}\) and \(\Omega_M=(\gamma_L-1)/(\gamma_L+1)\).  \(\square\)

\paragraph{Consequence.}
The most loaded assumption of the companion is the assignment of counting rules to the frames.  Deriving it from the frames (OP1 since v4.9.0-r63)
would strengthen three open predictions at once; a refutation of one of them points first to that rule.  If H15, H46 and H96 are all confirmed, this
supports the core more than one hit and less than three independent ones.  Assumption U is now explicit.

\subsection*{v4.10.0-r08 status in the three frames}
\label{subsec:v4100-r08-frames}

\begin{itemize}[leftmargin=2em]
\item M: r08.1 concerns the matter reading after the bang; all four predictions are read out in frame M.
\item T: enters r08.3 through \(\gamma_{\rm string}\) (counting \(j=\tfrac12\)) and the map from T to M; under G2 the tachyon reading lives before the bang.
\item P: no separate statement.
\end{itemize}
No frame-free statement is made.

\subsection*{v4.10.0-r08 register entries}
\label{subsec:v4100-r08-register}

\begin{center}
\small
\begin{tabular}{@{}>{\raggedright\arraybackslash}p{0.34\textwidth}>{\raggedright\arraybackslash}p{0.58\textwidth}@{}}
\toprule
statement & status after v4.10.0-r08\\
\midrule
second reading vs.\ photon weights of gravity & no conflict under G2 (r08.1)\\
equal weight of all fields after the bang (U) & tacit assumption, now stated\\
open predictions H15, H42, H46, H96 & all status reading; H15, H46, H96 on one \(\gamma_L\) core, H42 independent\\
most loaded assumption & counting rule per frame (OP1)\\
\bottomrule
\end{tabular}
\end{center}

\subsection*{v4.10.0-r08 ledger and guard update}
\label{subsec:v4100-r08-ledger-guard}

\[
R_{\rm chk}^{v4100\text{-}r08}=(1,1,0)
\]
for (no conflict of weights; dependence of the predictions traced; the counting rule not derived).

The layer spends no guard:
\[
R_{\rm glob}^{(v4100\text{-}r08,+)}=(0,1,1,1,1),
\qquad
\Gamma_{\rm guard}^{v4100\text{-}r08}=(1,1,1,1,1),
\qquad
\beta_4=1/256\ \text{unspent}.
\]
It does not imply exact finite \(C5\) zero closure, exact all-readout terminality,
\(C6/T2\), a global photon-clock/RH theorem, or global all-readout closure.
